\begin{answers}{ch:distillation}{蒸留と量子化：モデルを小さく速くする}
\item $s=0.4$、$z_0=1$、$q_i=(0,\,1,\,1,\,3)$、$\hat w_i=(-0.4,\,0.0,\,0.0,\,0.8)$、誤差は $(0,\,0.1,\,0.2,\,0)$ で、どれも $s/2=0.2$ 以下です。
スケールは $s=(0.8-(-0.4))/3=0.4$、ゼロ点は $z_0=\operatorname{round}(0.4/0.4)=1$ です。
$w_i/s=(-1,\,0.25,\,0.5,\,2)$ を丸めて $z_0$ を足すと $q_i$ になり、\eqref{eq:quant-deq}より $\hat w_i=0.4\,(q_i-1)$ です。
$w_3=0.2$ は $w_3/s=0.5$ で、目盛り $0.0$ と $0.4$ のちょうど真ん中に当たります。
ここでは \file{quantize.py} と同じ \code{torch.round}（ちょうど半分のときは偶数に丸める）を使い、$0.5$ を $0$ に丸めて $q_3=1$、$\hat w_3=0.0$ としました。
四捨五入なら $1$ に丸めて $q_3=2$、$\hat w_3=0.4$ です。誤差はどちらでも $0.2$ で、ちょうど上限の $s/2$ になります。
\item 勾配の大きさ（ノルム）は $0.421$ から $0.114$ になり、およそ $1/3.7$ です。
生徒のロジットはすべて等しいので、$\vq^\tau$ はどちらの温度でも $(1/3,\,1/3,\,1/3)$ です。
$\tau=1$ では $\vp^1=\softmax(2,1,0)=(0.665,\,0.245,\,0.090)$ で、勾配は $(\vq^1-\vp^1)/1=(-0.332,\,0.089,\,0.243)$ です。
$\tau=2$ では $\vp^2=\softmax(1,0.5,0)=(0.506,\,0.307,\,0.186)$ で、勾配は $(\vq^2-\vp^2)/2=(-0.087,\,0.013,\,0.073)$ です。
$1/\tau$ の係数だけなら $1/2$ のはずですが、$\vp^\tau$ が一様分布に近づいて差 $\vq^\tau-\vp^\tau$ も縮むので、\eqref{eq:distill-grad-hot}の $1/\tau^2$（$1/4$）に近い値になります。
\item 3\% のデータでは、$\tau$ を変えても最良の検証損失はほとんど変わらず、$\alpha$ を 1 にすると大きく下がりました（実測。学習は各条件 1 回で、最良はどれも 400 ステップ目）。
$\alpha=0.5$ のまま $\tau$ を 1、2、4 と変えると、$6.86$、$6.88$、$6.89$ です（1 と 2 は本文の値）。
$\tau=2$ で $\alpha=1$ にすると $6.50$、$\tau=1$、$\alpha=1$（\secref{sec:distill-llm}の Gemma 2 と同じ形）では $6.36$ でした。
3\% のデータでは、正解の項が丸暗記のほうへ生徒を引っぱるので、外したほうが良くなったと考えられます。
$\alpha=1$、$\tau=2$ で $\tau^2$ を外すと $6.49$ で、$6.50$ とほぼ同じです。
損失全体が $1/4$ 倍になるだけで、AdamW の 1 歩の大きさは、勾配の大小によらずほぼ学習率で決まるからです（\secref{sec:training-optim}）。
\item 差は、データを 30\% から 3\% へ減らすと $0.31$、$0.56$、$0.97$ と広がりますが、1\% では $0.52$ に縮みました（1\% と 30\% は実測、ほかは\tabref{tab:distill-results}）。
30\% は (a) $6.86$、(b) $6.55$ で、差は 100\% の $0.33$ とほぼ同じです。
1\%（13,632 トークン）の (a) は、50 ステップ目にもう $8.45$ で、学習前の $8.34$ より悪くなります。
50 ステップごとの記録では「最良」が学習前の値になるので、1\% は 5〜10 ステップごとに測り直しました。
(a) の最良は 10 ステップ目の $7.88$、(b) は 170 ステップ目の $7.36$ です（3\% の (a) も同じように測り直すと $7.83$ で、差は $0.95$ とほぼ変わりません）。
1\% では (b) も途中から悪化します。データが少なすぎると、教師の分布でも過学習を防ぎきれません。
横軸を割合（対数目盛り）、縦軸を差にしたグラフは、3\% に山のある形です。
\item 正答率は (a) とほぼ同じで、(a) を上回ったとは言えませんでした（実測。学習は各条件 1 回）。
教師の答えが正解と一致した対話だけを残すと、3,000 件のうち 3 桁までは 2,731 件、2 桁までは 2,730 件です。
これで (c) と同じように学習し、\tabref{tab:distill-drill}と同じ 1,000 問（学習データと同じ質問を 105 問と 190 問含みます）で測りました。
3 桁までは $7.9\%$（(a) $8.4\%$、選別なしの (c) $5.1\%$）、2 桁までは $24.3\%$（(a) $21.4\%$、(c) $20.2\%$）です。
同じ質問を除いた残りで数え直しても、$4.6\%$（(a) $4.8\%$）と $16.0\%$（(a) $13.7\%$）で、傾向は同じです。
2 桁の差 $2.9$ ポイントは評価の標準誤差（約 $1.3$ ポイント）の 2 倍ほどですが、3 桁では逆に低く、選別の効果とは言い切れません。
残した手本は正解と一字一句同じなので、この選別は「(a) から、教師が間違えた約 9\% を抜く」ことにしかなりません。
選別が役に立つのは、DeepSeek-R1 の思考過程のように、正解ラベルにない内容を教師が書く場合です。
\end{answers}
